\documentclass[12pt,a4paper]{article}

% --- Packages ---
\usepackage[utf8]{inputenc}
\usepackage[T1]{fontenc}
\usepackage[bahasa]{babel}
\usepackage{amsmath,amssymb}
\usepackage{graphicx}
\usepackage{booktabs}
\usepackage{geometry}
\usepackage{hyperref}
\usepackage{float}
\usepackage{caption}
\usepackage{subcaption}
\usepackage{listings}
\usepackage{xcolor}
\usepackage{enumitem}
\usepackage{tabularx}
\usepackage{cite}

\geometry{top=2.5cm,left=3cm,right=2.5cm,bottom=2.5cm}

\hypersetup{colorlinks=true,linkcolor=black,urlcolor=blue,citecolor=blue}

\lstset{
    basicstyle=\footnotesize\ttfamily,
    numbers=left,
    numberstyle=\tiny,
    frame=single,
    breaklines=true,
    backgroundcolor=\color{gray!10},
}

\title{\textbf{Metodologi Deteksi Deforestasi} \\
        Menggunakan \textit{Deep Learning} U-Net \\
        pada Citra Sentinel-2}
\author{Peneliti -- Deforest.id}
\date{}

\begin{document}
\maketitle
\thispagestyle{empty}

\begin{abstract}
Deforestasi merupakan salah satu permasalahan lingkungan yang kritis di Indonesia. 
Penelitian ini mengembangkan sistem deteksi deforestasi otomatis menggunakan arsitektur 
U-Net berbasis citra satelit Sentinel-2. Metodologi mencakup akuisisi data dari Google 
Earth Engine (GEE), pra-pemrosesan meliputi deteksi awan heuristik RGB dan NDVI, 
pembuatan \textit{pseudo-label} berbasis perubahan NDVI temporal, serta pelatihan 
model segmentasi semantik dengan 5 kanal masukan (RGB + NIR + NDVI). Model dilatih 
pada 11.500 sampel dari 4 scene, divalidasi pada 2.908 sampel, dan diuji pada 5.695 
sampel dari 7 scene Sentinel-2. Hasil pengujian menunjukkan IoU sebesar 0,7256, 
Dice 0,8410, Precision 0,8340, dan Recall 0,8480.
\end{abstract}

\section{Pendahuluan}
\label{sec:pendahuluan}

Indonesia memiliki hutan tropis terluas ketiga di dunia setelah Brasil dan 
Republik Demokratik Kongo. Namun, laju deforestasi di Indonesia masih tinggi 
akibat aktivitas illegal logging, perluasan perkebunan sawit, pertambangan, 
dan kebakaran hutan \cite{fao2020}. Deteksi dini dan pemantauan deforestasi 
secara periodik sangat diperlukan untuk mendukung upaya mitigasi dan penegakan 
hukum.

Citra satelit optik resolusi menengah seperti Sentinel-2 (ESA) menyediakan data 
multispektral dengan resolusi spasial 10 meter dan frekuensi temporal 5 hari, 
sehingga sangat cocok untuk pemantauan deforestasi \cite{drusch2012sentinel}. 
Pendekatan \textit{deep learning}, khususnya arsitektur U-Net \cite{ronneberger2015unet}, 
telah terbukti efektif untuk segmentasi semantik citra satelit.

Penelitian ini bertujuan untuk mengembangkan metodologi deteksi deforestasi 
otomatis yang mencakup seluruh pipeline dari akuisisi data hingga inferensi, 
dengan fokus pada:

\begin{enumerate}[nosep]
    \item Akuisisi dan pra-pemrosesan data Sentinel-2 multi-temporal
    \item Pembuatan \textit{pseudo-label} deforestasi berbasis perubahan NDVI
    \item Pelatihan model U-Net untuk segmentasi deforestasi
    \item Evaluasi dan visualisasi hasil prediksi
\end{enumerate}

\section{Data dan Materi}
\label{sec:data}

\subsection{Sumber Data}

Data yang digunakan berasal dari citra satelit Sentinel-2 Level-2A 
(\textit{Surface Reflectance}) yang diakses melalui Google Earth Engine (GEE). 
Citra dipilih dari 7 scene yang mencakup wilayah hutan di Indonesia dengan 
kombinasi dua periode waktu: \textit{baseline} (kondisi hutan awal) dan 
\textit{post-deforestation} (kondisi setelah deforestasi).

\subsubsection{Kanal Spektral}

Empat kanal spektral Sentinel-2 digunakan:

\begin{itemize}[nosep]
    \item \textbf{B2 (Blue)}: 490 nm, resolusi 10 m
    \item \textbf{B3 (Green)}: 560 nm, resolusi 10 m
    \item \textbf{B4 (Red)}: 665 nm, resolusi 10 m
    \item \textbf{B8 (NIR)}: 842 nm, resolusi 10 m
\end{itemize}

Selain itu, \textit{Normalized Difference Vegetation Index} (NDVI) dihitung 
sebagai kanal tambahan:

\begin{equation}
    \label{eq:ndvi}
    \text{NDVI} = \frac{\text{NIR} - \text{Red}}{\text{NIR} + \text{Red}}
\end{equation}

\subsection{Studi Area}

Tujuh scene Sentinel-2 digunakan dalam penelitian ini. Setiap scene mencakup 
area sekitar $109,8 \times 109,8$ km$^2$ (10.980 $\times$ 10.980 piksel pada 
resolusi 10 m), atau sekitar 1,2 juta hektar per scene. Scene-scene tersebut 
dipliih berdasarkan ketersediaan data multi-temporal dan bukti deforestasi yang 
terverifikasi.

\begin{table}[H]
\centering
\caption{Sebaran scene berdasarkan split dataset}
\label{tab:scenes}
\begin{tabular}{lrc}
\toprule
\textbf{Split} & \textbf{Jumlah Scene} & \textbf{Scene} \\
\midrule
Train & 4 & Beragam \\
Val   & 1 & Mukomuko tahun 2019 \\
Test  & 2 & Beragam \\
\bottomrule
\end{tabular}
\end{table}

\subsection{Pembagian Dataset}

Dataset dibagi dengan stratifikasi berdasarkan scene untuk menghindari 
\textit{data leakage} antar split:

\begin{itemize}[nosep]
    \item \textbf{Train}: 11.500 sampel (57,2\%) dari 4 scene
    \item \textbf{Validation}: 2.908 sampel (14,5\%) dari 1 scene
    \item \textbf{Test}: 5.695 sampel (28,3\%) dari 2 scene
\end{itemize}

Total: \textbf{20.103 sampel} dari 7 scene unik. Masing-masing sampel berupa 
\textit{chip} berukuran $64 \times 64$ piksel setara dengan $640 \times 640$ m 
di lapangan.

\section{Metodologi}
\label{sec:metodologi}

\subsection{Gambaran Umum Pipeline}

Pipeline pemrosesan terdiri dari enam tahap utama:

\begin{enumerate}[nosep]
    \item Akuisisi data Sentinel-2 dari GEE
    \item \textit{Tiling} citra menjadi tile $512 \times 512$ piksel
    \item \textit{Chipping} menjadi chip $64 \times 64$ piksel
    \item Deteksi awan dan pembuatan masker
    \item Pembuatan \textit{pseudo-label} deforestasi berbasis NDVI
    \item Pelatihan model U-Net
\end{enumerate}

\subsection{Akuisisi dan Pra-pemrosesan Data}

\subsubsection{Tiling}

Setiap scene Sentinel-2 berukuran 10.980 $\times$ 10.980 piksel di-\textit{tile}
menjadi tile berukuran $512 \times 512$ piksel. Proses tiling dilakukan dengan 
\textit{sliding window} tanpa tumpang tindih (\textit{overlap}), menghasilkan 
sekitar 440 tile per scene.

\subsubsection{Chipping}

Tile $512 \times 512$ piksel kemudian di-\textit{chip} menjadi $64 \times 64$ 
piksel dengan \textit{stride} 32 piksel (50\% overlap). \textit{Overlap} ini 
bertujuan untuk:

\begin{itemize}[nosep]
    \item Meningkatkan jumlah sampel pelatihan
    \item Mengurangi risiko deforestasi terpotong di tepi chip
    \item Memberikan konteks spasial yang memadai untuk input U-Net
\end{itemize}

\subsection{Deteksi Awan}

Kanal QA60 (\textit{cloud mask band}) pada data Sentinel-2 yang digunakan 
seluruhnya bernilai nol (tidak berfungsi), sehingga deteksi awan dilakukan 
menggunakan heuristik berbasis RGB dan NDVI.

Suatu piksel diklasifikasikan sebagai awan jika memenuhi tiga kondisi:

\begin{equation}
    \label{eq:cloud}
    \text{Awan}(p) =
    \begin{cases}
        1, & \text{jika } R > T_{bright} \land G > T_{bright} \land B > T_{bright} \\
           & \land \text{NDVI} < T_{ndvi} \land B \geq G \geq R \\
        0, & \text{lainnya}
    \end{cases}
\end{equation}

dengan ambang batas $T_{bright} = 180$ (dari rentang 0--255) dan $T_{ndvi} = 0,10$.

Alasan heuristik ini:

\begin{itemize}[nosep]
    \item Awan memantulkan seluruh spektrum tampak dengan intensitas tinggi
    \item NDVI awan mendekati nol karena tidak ada vegetasi
    \item Awan cenderung memiliki \textit{blue shift} (B $>$ G $>$ R)
\end{itemize}

Piksel awan diberi label \texttt{IGNORE\_LABEL = 255} dan tidak diikutsertakan 
dalam perhitungan \textit{loss} selama pelatihan.

\subsection{Pembuatan Pseudo-label Deforestasi}

\label{sec:label}

\textit{Pseudo-label} deforestasi dibuat menggunakan perubahan NDVI temporal 
antara citra \textit{baseline} (hutan utuh) dan \textit{post-deforestation} 
(dugaan deforestasi).

\subsubsection{NDVI Change Detection}

Perubahan NDVI dihitung sebagai:

\begin{equation}
    \label{eq:ndvi_change}
    \Delta\text{NDVI} = \text{NDVI}_{t2} - \text{NDVI}_{t1}
\end{equation}

dengan $t_1$ = \textit{baseline} dan $t_2$ = \textit{post-deforestation}. 
Suatu piksel diklasifikasikan sebagai deforestasi jika:

\begin{equation}
    \label{eq:deforest}
    \text{Deforestasi}(p) =
    \begin{cases}
        1, & \text{jika } \Delta\text{NDVI} < T_{ndvi\_change} \\
        0, & \text{lainnya}
    \end{cases}
\end{equation}

dengan ambang batas $T_{ndvi\_change} = -0,15$. Nilai negatif menunjukkan 
penurunan NDVI yang signifikan, mengindikasikan hilangnya vegetasi.

\subsubsection{Post-processing Morfologi}

Masker deforestasi hasil deteksi perubahan NDVI menjalani dua tahap 
post-processing morfologi:

\begin{enumerate}[nosep]
    \item \textbf{Morphological opening} (erosion $\rightarrow$ dilation) 
          dengan kernel ellipse $5 \times 5$ untuk menghilangkan \textit{salt noise}
    \item \textbf{Morphological closing} (dilation $\rightarrow$ erosion) 
          dengan kernel yang sama untuk mengisi lubang kecil pada area deforestasi
    \item \textbf{Connected component filtering}: komponen dengan luas 
          $< 64$ piksel ($< 0,64$ ha) dihapus
\end{enumerate}

\subsection{Arsitektur U-Net}

\subsubsection{Struktur Model}

Model yang digunakan adalah U-Net \cite{ronneberger2015unet} dengan 
5 kanal masukan dan 2 kanal keluaran. Arsitektur terdiri dari:

\begin{itemize}[nosep]
    \item \textbf{Encoder}: 4 tahap \textit{down-sampling}, masing-masing 
          terdiri dari 2 lapisan konvolusi $3 \times 3$ + Batch Normalization 
          + ReLU, diikuti \textit{max pooling} $2 \times 2$
    \item \textbf{Bottleneck}: 512 filter konvolusi
    \item \textbf{Decoder}: 4 tahap \textit{up-sampling} dengan 
          \textit{transposed convolution} $2 \times 2$ + \textit{skip connection} 
          dari encoder
    \item \textbf{Output layer}: konvolusi $1 \times 1$ untuk memetakan ke 
          2 kelas (forest/deforest)
\end{itemize}

\begin{table}[H]
\centering
\caption{Detail arsitektur U-Net}
\label{tab:unet_arch}
\begin{tabular}{lccc}
\toprule
\textbf{Tahap} & \textbf{Input Channels} & \textbf{Output Channels} & \textbf{Ukuran} \\
\midrule
Input & 5 & -- & $64 \times 64$ \\
Encoder 1 & 5 & 32 & $64 \times 64$ \\
Encoder 2 & 32 & 64 & $32 \times 32$ \\
Encoder 3 & 64 & 128 & $16 \times 16$ \\
Encoder 4 & 128 & 256 & $8 \times 8$ \\
Bottleneck & 256 & 512 & $4 \times 4$ \\
Decoder 4 & 512+256 & 256 & $8 \times 8$ \\
Decoder 3 & 256+128 & 128 & $16 \times 16$ \\
Decoder 2 & 128+64 & 64 & $32 \times 32$ \\
Decoder 1 & 64+32 & 32 & $64 \times 64$ \\
Output & 32 & 2 & $64 \times 64$ \\
\bottomrule
\end{tabular}
\end{table}

Total parameter model: $\sim$7,5 juta.

\subsection{Pelatihan Model}

\subsubsection{Fungsi Loss}

Fungsi loss yang digunakan adalah \textit{Dice Loss} yang dirancang untuk 
menangani ketidakseimbangan kelas (\textit{class imbalance}):

\begin{equation}
    \label{eq:dice_loss}
    \mathcal{L}_{Dice} = 1 - \frac{1}{C} \sum_{c=1}^{C} 
    \frac{2 \sum_{i} p_{i,c} \cdot g_{i,c} + \epsilon}{\sum_{i} p_{i,c} + \sum_{i} g_{i,c} + \epsilon}
\end{equation}

dengan $C=2$ (jumlah kelas), $p_{i,c}$ adalah probabilitas prediksi 
(setelah softmax), $g_{i,c}$ adalah \textit{ground truth} one-hot, 
dan $\epsilon = 1,0$ adalah \textit{smoothing factor}.

Piksel dengan label 255 (\textit{cloud regions}) diabaikan dalam perhitungan loss.

\subsubsection{Hiperparameter}

\begin{table}[H]
\centering
\caption{Hiperparameter pelatihan}
\label{tab:hyperparams}
\begin{tabular}{lc}
\toprule
\textbf{Parameter} & \textbf{Nilai} \\
\midrule
Jumlah epoch & 50 \\
Batch size & 32 \\
Learning rate awal & $1 \times 10^{-3}$ \\
Scheduler & CosineAnnealingLR ($T_{max}=50$) \\
Optimizer & AdamW ($\beta_1=0,9$, $\beta_2=0,999$) \\
Weight decay & $1 \times 10^{-4}$ \\
Input size & $5 \times 64 \times 64$ \\
num\_workers & 0 \\
\bottomrule
\end{tabular}
\end{table}

\subsubsection{Normalisasi Input}

Setiap kanal input dinormalisasi sebagai berikut:

\begin{itemize}[nosep]
    \item \textbf{RGB}: $\text{piksel} / 255,0$ (rentang 0--1)
    \item \textbf{NIR}: $\text{piksel} / 10000,0$ karena nilai \textit{surface 
          reflectance} Sentinel-2 berkisar 0--10000
    \item \textbf{NDVI}: langsung digunakan (rentang -1--1)
\end{itemize}

Piksel dengan nilai NaN (\textit{Not a Number}) akibat artefak data diganti 
dengan 0 melalui fungsi \texttt{np.nan\_to\_num}.

\subsubsection{Perangkat Keras}

\begin{itemize}[nosep]
    \item \textbf{GPU}: NVIDIA RTX 3060 (12 GB VRAM)
    \item \textbf{Sistem Operasi}: Ubuntu 26.04 (WSL2 di Windows 11)
    \item \textbf{CPU}: Intel (12 core)
    \item \textbf{RAM}: 32 GB
\end{itemize}

Kecepatan pelatihan: $\sim$2,7 menit per epoch (360 iterasi train, 
91 iterasi validasi per epoch).

\subsection{Evaluasi}

\subsubsection{Metrik}

Empat metrik evaluasi digunakan:

\begin{itemize}[nosep]
    \item \textbf{Intersection over Union (IoU)}:
    \begin{equation}
        \text{IoU} = \frac{TP}{TP + FP + FN}
    \end{equation}
    
    \item \textbf{Dice Coefficient}:
    \begin{equation}
        \text{Dice} = \frac{2 \times TP}{2 \times TP + FP + FN}
    \end{equation}
    
    \item \textbf{Precision}:
    \begin{equation}
        \text{Precision} = \frac{TP}{TP + FP}
    \end{equation}
    
    \item \textbf{Recall}:
    \begin{equation}
        \text{Recall} = \frac{TP}{TP + FN}
    \end{equation}
\end{itemize}

dengan TP = \textit{True Positive}, FP = \textit{False Positive}, 
FN = \textit{False Negative}.

\subsubsection{Prosedur Validasi}

Validasi dilakukan per epoch pada validation set (1 scene, 2.908 sampel). 
Model dengan IoU validasi tertinggi disimpan sebagai \texttt{best.pth}. 
Setelah seluruh epoch selesai, model final disimpan sebagai \texttt{final.pth}.

Evaluasi akhir dilakukan pada test set (2 scene, 5.695 sampel) menggunakan 
model terbaik.

\section{Hasil dan Pembahasan}
\label{sec:hasil}

\subsection{Progres Pelatihan}

\begin{table}[H]
\centering
\caption{Progres pelatihan per epoch}
\label{tab:training_progress}
\begin{tabular}{cccc}
\toprule
\textbf{Epoch} & \textbf{Train Loss} & \textbf{Val Loss} & \textbf{Val IoU} \\
\midrule
1 & 0,1974 & 0,2753 & 0,3788 \\
5 & 0,1775 & 0,2773 & 0,3726 \\
\textbf{9} & \textbf{0,1759} & \textbf{0,2765} & \textbf{0,3831*} \\
15 & 0,1704 & 0,2792 & 0,3628 \\
25 & 0,1655 & 0,2928 & 0,3653 \\
37 & 0,1598 & 0,2859 & 0,3649 \\
50 & 0,1555 & 0,2865 & 0,3616 \\
\bottomrule
\multicolumn{4}{l}{\small *Model terbaik (\texttt{best.pth}) disimpan pada epoch ini}
\end{tabular}
\end{table}

\textit{Train loss} menurun secara konsisten dari 0,1974 ke 0,1555, menunjukkan 
model terus belajar. \textit{Val IoU} tertinggi dicapai pada epoch 9 sebesar 
0,3831, setelah itu cenderung stagnan di kisaran 0,36.

\subsection{Hasil Test Set}

\begin{table}[H]
\centering
\caption{Metrik evaluasi pada test set}
\label{tab:test_metrics}
\begin{tabular}{lc}
\toprule
\textbf{Metrik} & \textbf{Nilai} \\
\midrule
IoU & 0,7256 \\
Dice & 0,8410 \\
Precision & 0,8340 \\
Recall & 0,8480 \\
\bottomrule
\end{tabular}
\end{table}

\subsection{Analisis Hasil}

Terdapat perbedaan signifikan antara IoU validasi (0,38) dan IoU test (0,73). 
Hal ini disebabkan oleh:

\begin{enumerate}[nosep]
    \item \textbf{Stratifikasi berdasarkan scene}: Validation set hanya terdiri 
          dari 1 scene (Mukomuko tahun 2019) yang memiliki karakteristik medan 
          berbeda dengan scene-scene lain
    \item \textbf{Karakteristik scene}: Scene validasi mungkin memiliki kondisi 
          tutupan lahan, topografi, atau pola deforestasi yang lebih sulit 
          dideteksi
    \item \textbf{Distribusi IoU bimodal}: Distribusi IoU pada test set 
          cenderung bimodal — banyak sampel dengan IoU sangat tinggi ($>0,9$) 
          dan sangat rendah ($<0,3$)
\end{enumerate}

\subsection{Analisis Error}

Tiga mode kegagalan utama teridentifikasi:

\begin{enumerate}[nosep]
    \item \textbf{Miss (False Negative dominan)}: Model gagal mendeteksi area 
          deforestasi yang kecil atau deforestasi halus. Ini umum terjadi pada 
          deforestasi dengan ukuran di bawah 5 piksel atau perubahan NDVI 
          yang tidak terlalu drastis.
    \item \textbf{False Positive}: Model salah mengklasifikasikan area non-hutan 
          sebagai deforestasi. Area yang sering menjadi FP meliputi:
          \begin{itemize}[nosep]
              \item Badan sungai kering atau tepi sungai
              \item Area persawahan yang tidak berpohon
              \item Lahan terbuka
              \item Tepi awan (\textit{cloud edge}) yang lolos deteksi awan
          \end{itemize}
    \item \textbf{Boundary Error}: Model kurang presisi pada batas antara area 
          hutan dan deforestasi. Ini adalah keterbatasan umum pada model 
          segmentasi dengan resolusi piksel 10 meter.
\end{enumerate}

\section{Kesimpulan}
\label{sec:kesimpulan}

Penelitian ini berhasil mengembangkan pipeline deteksi deforestasi otomatis 
berbasis U-Net pada citra Sentinel-2. Model mencapai IoU 0,7256 dan Dice 0,8410 
pada test set, menunjukkan kelayakan pendekatan ini untuk deteksi deforestasi 
skala operasional.

Kontribusi utama penelitian ini meliputi:

\begin{enumerate}[nosep]
    \item Pipeline lengkap dari akuisisi GEE hingga inferensi
    \item Metode \textit{pseudo-labeling} berbasis perubahan NDVI temporal 
          yang memungkinkan pembuatan dataset tanpa anotasi manual
    \item Heuristik deteksi awan berbasis RGB dan NDVI sebagai solusi atas 
          kanal QA60 yang tidak berfungsi
    \item Analisis komprehensif mode kegagalan model pada berbagai kondisi 
          deforestasi
\end{enumerate}

Untuk pengembangan selanjutnya, beberapa arah penelitian yang potensial meliputi:

\begin{itemize}[nosep]
    \item \textbf{Augmentasi data}: Rotasi, flipping, dan \textit{color jitter} 
          untuk meningkatkan generalisasi
    \item \textbf{Ekstraksi fitur temporal}: Menggunakan rangkaian multi-temporal 
          daripada pasangan baseline-deforest
    \item \textbf{Arsitektur lanjutan}: Swin U-Net atau Transformer-based 
          segmentation untuk menangkap konteks global
    \item \textbf{Fusion dengan data SAR}: Sentinel-1 (radar) untuk menembus 
          tutupan awan
\end{itemize}

\bibliographystyle{ieeetr}
\begin{thebibliography}{3}

\bibitem{fao2020}
FAO,
\emph{Global Forest Resources Assessment 2020: Main report},
Rome, 2020.

\bibitem{drusch2012sentinel}
M.~Drusch et al.,
``Sentinel-2: ESA's Optical High-Resolution Mission for GMES Operational Services,''
\emph{Remote Sensing of Environment}, vol.~120, pp.~25--36, 2012.

\bibitem{ronneberger2015unet}
O.~Ronneberger, P.~Fischer, and T.~Brox,
``U-Net: Convolutional Networks for Biomedical Image Segmentation,''
in \emph{Proc. MICCAI}, 2015, pp.~234--241.

\end{thebibliography}

\end{document}
